%% FullCourse_10Lecs -- Lecture 9, Part 9.2a
%% Assembled by tools/lec09_assemble.py from reorg_manifest.yaml: Phase A of
%% Lec09_Reorg_Plan_2026-09-12.md as amended by Lec09_Reorg_Plan_Review_2026-09-12.md.
%% Each "% ---- <ID> · <TAG> · TIER" marker names a frame's source deck and line;
%% keep the markers until Phase B is done (lec09_assemble.py check reads them).
%% Owns: harness/LLM/files, the JSON tool call, the privacy boundary, session-as-list, the twenty-line loop, ReAct
%% Owns: {'the GDP-by-party worked trace (recorded transcript': 'labs/Lec09_Agent_Lab/transcripts/)'}
%% Defers: install, files, commands, permissions, tokens/context mechanics -> T2b; dials, ceilings, dispatch -> T3b

%% Shared preamble for the FullCourse_10Lecs topic decks.
%%
%% Each topic .tex \input's this file, then defines its own \title, \date
%% override (if any), \graphicspath, and \begin{document}.
%%
%% Reconstructed verbatim from the inlined copy preserved in
%% Lec10_Case_Studies/Lec10_T8_Case6_DeepRamsey.tex (lines 23-190), which is
%% explicitly annotated there as "from ../shared_preamble.tex, copied verbatim".
%% Base preamble matches MiniCourse_8hr/Slides/shared_preamble.tex; this version
%% additionally provides \sectiondivider, the conceptbox/cautionbox/examplebox/
%% takeawaybox tcolorboxes, and \compacttables used across the split decks.

\documentclass[10pt,english,aspectratio=169]{beamer}

\usepackage{ae,aecompl}
\usepackage[T1]{fontenc}
\usepackage{textcomp}
\usepackage[utf8]{inputenc}
\usepackage{babel}
\usepackage{datetime2}

\setcounter{secnumdepth}{3}
\setcounter{tocdepth}{3}

\usepackage{amsmath, amssymb, amsthm, graphicx}
\usepackage{booktabs, multirow}
\usepackage{tikz}
\usepackage{pgfplots}
\pgfplotsset{compat=1.16}
\usepackage[most]{tcolorbox}
\tcbuselibrary{skins, breakable, theorems}
\usepackage{listings}
\usepackage{xcolor}

\lstset{
  basicstyle=\ttfamily\small,
  keywordstyle=\color{blue!80!black}\bfseries,
  commentstyle=\color{green!50!black}\itshape,
  stringstyle=\color{orange!80!black},
  backgroundcolor=\color{gray!8},
  frame=single,
  rulecolor=\color{gray!40},
  breaklines=true,
  showstringspaces=false,
  numbers=none,
}

\lstdefinestyle{pythonstyle}{
    language=Python,
    basicstyle=\ttfamily\footnotesize,
    keywordstyle=\color{blue!80!black}\bfseries,
    commentstyle=\color{green!50!black}\itshape,
    stringstyle=\color{orange!80!black},
    frame=single,
    breaklines=true,
    showstringspaces=false,
    numbers=left,
    numbersep=5pt,
    numberstyle=\tiny\color{gray},
    backgroundcolor=\color{gray!8},
    rulecolor=\color{gray!40}
}

\ifx\hypersetup\undefined
  \AtBeginDocument{%
    \hypersetup{bookmarks=false,breaklinks=true,pdfborder={0 0 0},colorlinks=true,
      linkcolor=DarkText,urlcolor=RedTitle}%
  }
\else
  \hypersetup{bookmarks=false,breaklinks=true,pdfborder={0 0 0},colorlinks=true,
    linkcolor=DarkText,urlcolor=RedTitle}%
\fi

\makeatletter
\newcommand\makebeamertitle{%
  \begin{frame}[plain]
    \vspace*{1.0cm}
    \centering
    {\usebeamerfont{title}\usebeamercolor[fg]{title}\inserttitle\par}
    \vspace{1.0cm}
    {\usebeamerfont{author}\usebeamercolor[fg]{author}\insertauthor\par}
    \vspace{0.2cm}
    {\usebeamerfont{date}\usebeamercolor[fg]{date}\insertdate\par}
  \end{frame}}%
\AtBeginDocument{%
  \let\origtableofcontents=\tableofcontents
  \def\tableofcontents{\@ifnextchar[{\origtableofcontents}{\gobbletableofcontents}}
  \def\gobbletableofcontents#1{\origtableofcontents}%
}

\usetheme{metropolis}
\usefonttheme{professionalfonts}

\definecolor{LightGray}{RGB}{230,230,230}
\definecolor{RedTitle}{RGB}{0,0,200}
\definecolor{VeryLightGray}{RGB}{245,245,245}
\definecolor{DarkText}{RGB}{0,0,0}

\setbeamercolor{background canvas}{bg=white}
\setbeamercolor{title}{fg=RedTitle}
\setbeamercolor{author}{fg=DarkText}
\setbeamercolor{date}{fg=DarkText}
\setbeamercolor{frametitle}{bg=VeryLightGray,fg=DarkText}
\setbeamercolor{normal text}{fg=DarkText}
\setbeamerfont{title}{size=\large}

\setbeamersize{text margin left=5mm, text margin right=5mm}
\addtobeamertemplate{frametitle}{}{\vspace{-0.5em}}
\newcommand{\reducevspace}{\vspace{-0.5cm}}

% Shared section divider and box styles for split topic decks.
\newcommand{\sectiondivider}[2]{%
\begin{frame}[plain,noframenumbering]
\begin{center}
\vfill
{\Large\textbf{#1}\par}
\vspace{0.35cm}
{\normalsize #2\par}
\vfill
\end{center}
\end{frame}
}

\newtcolorbox{conceptbox}[1][]{
  colback=blue!3,
  colframe=RedTitle!55,
  boxrule=0.35mm,
  arc=1mm,
  left=2mm,
  right=2mm,
  top=1mm,
  bottom=1mm,
  #1
}
\newtcolorbox{cautionbox}[1][]{
  colback=red!3,
  colframe=red!50!black,
  boxrule=0.35mm,
  arc=1mm,
  left=2mm,
  right=2mm,
  top=1mm,
  bottom=1mm,
  #1
}
\newtcolorbox{examplebox}[1][]{
  colback=green!3,
  colframe=green!45!black,
  boxrule=0.35mm,
  arc=1mm,
  left=2mm,
  right=2mm,
  top=1mm,
  bottom=1mm,
  #1
}
\newtcolorbox{takeawaybox}[1][]{
  colback=VeryLightGray,
  colframe=RedTitle!55,
  boxrule=0.35mm,
  arc=1mm,
  left=2mm,
  right=2mm,
  top=1mm,
  bottom=1mm,
  #1
}
\newcommand{\compacttables}{\renewcommand{\arraystretch}{1.15}}

\setbeamertemplate{navigation symbols}{}
\setbeamertemplate{footline}{%
  \hfill%
  \usebeamercolor[fg]{page number in head/foot}%
  \usebeamerfont{page number in head/foot}%
  \insertframenumber\,/\,\inserttotalframenumber\kern1em\vskip2pt%
}

\makeatother

\author{Zhigang Feng}
% "date with time": compile timestamp so each build is identifiable.
\date{\today\ \DTMcurrenttime}


% Suppress redundant auto-generated metropolis section page; \sectiondivider frames are the dividers we keep.
\metroset{sectionpage=none}

\graphicspath{{./pic/}{./figures/}{../}}

\usetikzlibrary{positioning,arrows.meta,shapes,calc,fit,backgrounds}

%% Lec09 shared MAP slide, v2 (September 2026 reorganization): five parts,
%% eleven decks.  v1 (the July "five acts" map) is archived as
%% _archive/five_acts_2026-07/lec09_map_v1.tex.
%%
%% Input this file AFTER ../shared_preamble.tex (needs RedTitle, VeryLightGray,
%% DarkText) and call \LecNineMap{part}{sub} inside a frame:
%%   part in {1..5} highlights that part ("you are here") and, when sub names a
%%   sub-deck letter, that deck's chip -- \LecNineMap{2}{b} is T2b, \LecNineMap{4}{}
%%   is T4;  part = 0 renders the closing reprise with every part complete (the
%%   "Your Job Description" frame in T5d).
%% Each part carries its student question and its economic twin (the delegation-
%% economics thread of Lec09_Reorg_Plan_Review_2026-09-12.md, A3).
%% Filename deliberately does not match the build glob Lec*_T*.tex, so the build
%% script never tries to compile it standalone.

% Highlight chip #1 when the requested deck key equals #2 (e.g. 2b).
\newcommand{\lecmapchip}[2]{%
  \edef\lecmaptmp{#2}%
  \ifx\lecmapkey\lecmaptmp\tikzset{#1/.style={lecchipcur}}\fi}

\newcommand{\LecNineMap}[2]{%
% TikZ option lists are comma-split before TeX conditionals run, so every
% \ifnum/\ifx is resolved here, into named styles, before the picture starts.
\edef\lecmapkey{#1#2}%
\tikzset{lecparta/.style={}, lecpartb/.style={}, lecpartc/.style={},
         lecpartd/.style={}, lecparte/.style={},
         chipTone/.style={}, chipTtwoa/.style={}, chipTtwob/.style={},
         chipTthreea/.style={}, chipTthreeb/.style={}, chipTfour/.style={},
         chipTfivea/.style={}, chipTfiveb/.style={}, chipTfivec/.style={},
         chipTfived/.style={}}%
\ifnum#1=0
  \tikzset{lecparta/.style={lecpartdone}, lecpartb/.style={lecpartdone},
           lecpartc/.style={lecpartdone}, lecpartd/.style={lecpartdone},
           lecparte/.style={lecpartdone}, lecchip/.style={lecchipbase, lecchipdone}}%
\else
  \tikzset{lecchip/.style={lecchipbase}}%
  \ifnum#1=1 \tikzset{lecparta/.style={lecpartcur}}\fi
  \ifnum#1=2 \tikzset{lecpartb/.style={lecpartcur}}\fi
  \ifnum#1=3 \tikzset{lecpartc/.style={lecpartcur}}\fi
  \ifnum#1=4 \tikzset{lecpartd/.style={lecpartcur}}\fi
  \ifnum#1=5 \tikzset{lecparte/.style={lecpartcur}}\fi
  \lecmapchip{chipTone}{1}\lecmapchip{chipTtwoa}{2a}\lecmapchip{chipTtwob}{2b}%
  \lecmapchip{chipTthreea}{3a}\lecmapchip{chipTthreeb}{3b}\lecmapchip{chipTfour}{4}%
  \lecmapchip{chipTfivea}{5a}\lecmapchip{chipTfiveb}{5b}%
  \lecmapchip{chipTfivec}{5c}\lecmapchip{chipTfived}{5d}%
\fi
\begin{center}
\begin{tikzpicture}[
    part/.style={rectangle, rounded corners, draw=black!60, fill=VeryLightGray,
        minimum width=2.62cm, minimum height=0.72cm, text width=2.5cm,
        align=center, font=\scriptsize, text=DarkText},
    lecpartcur/.style={draw=RedTitle, very thick, fill=orange!20},
    lecpartdone/.style={draw=green!45!black, thick, fill=green!10},
    lecchipbase/.style={rectangle, rounded corners=1pt, draw=black!30, fill=white,
        inner xsep=1.5pt, inner ysep=1pt, font=\tiny, text=black!70},
    lecchipcur/.style={draw=RedTitle, fill=RedTitle!12, text=RedTitle, font=\tiny\bfseries},
    lecchipdone/.style={draw=green!45!black, fill=green!8},
    q/.style={font=\tiny\itshape, text width=2.7cm, align=center, text=black!75},
    twin/.style={font=\tiny, text width=2.7cm, align=center, text=blue!40!black},
    sat/.style={rectangle, rounded corners, draw=black!45, dashed, fill=white,
        text width=5.6cm, align=center, font=\tiny, text=black!70, minimum height=0.42cm},
    barlab/.style={font=\tiny\bfseries, text=black!60},
    arr/.style={-{Stealth[length=2mm]}, thick, gray!60}
]
    % --- five part boxes ---
    \node[part, lecparta] (p1) at (0,0)     {\textbf{9.1 SEE}};
    \node[part, lecpartb] (p2) at (2.85,0)  {\textbf{9.2 UNDER THE HOOD}};
    \node[part, lecpartc] (p3) at (5.7,0)   {\textbf{9.3 BUILD}};
    \node[part, lecpartd] (p4) at (8.55,0)  {\textbf{9.4 DESIGN \& VALIDATE}};
    \node[part, lecparte] (p5) at (11.4,0)  {\textbf{9.5 RUN \& GOVERN}};
    \draw[arr] (p1) -- (p2);
    \draw[arr] (p2) -- (p3);
    \draw[arr] (p3) -- (p4);
    \draw[arr] (p4) -- (p5);
    % --- sub-deck chips ---
    \node[lecchip, chipTone]    at (0,-0.6)     {T1};
    \node[lecchip, chipTtwoa]   at (2.55,-0.6)  {T2a};
    \node[lecchip, chipTtwob]   at (3.15,-0.6)  {T2b};
    \node[lecchip, chipTthreea] at (5.4,-0.6)   {T3a};
    \node[lecchip, chipTthreeb] at (6.0,-0.6)   {T3b};
    \node[lecchip, chipTfour]   at (8.55,-0.6)  {T4};
    \node[lecchip, chipTfivea]  at (10.5,-0.6)  {T5a};
    \node[lecchip, chipTfiveb]  at (11.1,-0.6)  {T5b};
    \node[lecchip, chipTfivec]  at (11.7,-0.6)  {T5c};
    \node[lecchip, chipTfived]  at (12.3,-0.6)  {T5d};
    % --- the student question each part answers ---
    \node[q] at (0,-1.08)     {what changed,\\ and why care?};
    \node[q] at (2.85,-1.08)  {how does it run?\\ how do I start?};
    \node[q] at (5.7,-1.08)   {how do I build one,\\ and call it?};
    \node[q] at (8.55,-1.08)  {can I design one\\ and prove it works?};
    \node[q] at (11.4,-1.08)  {run a project, catch\\ mistakes, stay safe};
    % --- its economic twin ---
    \node[twin] at (0,-1.62)     {generation got cheap;\\ verification is scarce};
    \node[twin] at (2.85,-1.62)  {what am I paying for?\\ tokens, scarce context};
    \node[twin] at (5.7,-1.62)   {dividing the labour:\\ specialists, boundaries};
    \node[twin] at (8.55,-1.62)  {write the contract,\\ audit the work};
    \node[twin] at (11.4,-1.62)  {hidden action: gates,\\ monitoring, priors};
    % --- "you are here" marker above the current part ---
    \ifnum#1=1 \node[font=\scriptsize\bfseries, text=RedTitle] at (0,0.62)
        {you are here\ $\blacktriangledown$};\fi
    \ifnum#1=2 \node[font=\scriptsize\bfseries, text=RedTitle] at (2.85,0.62)
        {you are here\ $\blacktriangledown$};\fi
    \ifnum#1=3 \node[font=\scriptsize\bfseries, text=RedTitle] at (5.7,0.62)
        {you are here\ $\blacktriangledown$};\fi
    \ifnum#1=4 \node[font=\scriptsize\bfseries, text=RedTitle] at (8.55,0.62)
        {you are here\ $\blacktriangledown$};\fi
    \ifnum#1=5 \node[font=\scriptsize\bfseries, text=RedTitle] at (11.4,0.62)
        {you are here\ $\blacktriangledown$};\fi
    \ifnum#1=0 \node[font=\scriptsize\bfseries, text=green!45!black] at (5.7,0.62)
        {all five parts complete};\fi
    % --- bottom band: the three questions of the lecture ---
    \fill[blue!10]   (-1.3,-2.37) rectangle (4.15,-2.07);
    \fill[green!10]  (4.4,-2.37)  rectangle (9.85,-2.07);
    \fill[orange!15] (10.1,-2.37) rectangle (12.7,-2.07);
    \node[barlab] at (1.425,-2.22)  {HOW IT WORKS};
    \node[barlab] at (7.125,-2.22)  {HOW TO BUILD \& USE IT};
    \node[barlab] at (11.4,-2.22)   {YOUR ROLE};
    % --- dashed satellites ---
    \node[sat] at (1.9,-2.8)
        {\textbf{Appendix TA} --- the dated landscape and setup details, read on demand};
    \node[sat] at (9.9,-2.8)
        {\textbf{Lab} Parts A--B, Design Lab scenarios $\to$ \textbf{Lec10} cases (same morning)};
\end{tikzpicture}
\end{center}
}


\title[AI for Econ Research]{AI for Economic Research: Dynamic Models, Language, and Agents\\
Lecture 9.2a: How an Agent Actually Runs --- Harness, Session, Loop}

\begin{document}

\makebeamertitle

% ---- T2-F01 · MERGE-PENDING(T2-F01+T2b-F01) · TIER: CORE · from T2:31
% TODO(Phase B): One map frame: T2-F01's lead + \LecNineMap{2}{a}; fold T2b-F01's "corrects three beliefs that are widespread and wrong" into one line pointing to T3b's corrections table; then delete T2b-F01.
% TODO(Phase B): New route line: three components and one JSON line -> the session is a list -> the GDP trace. Add the 9.2 economic twin.
\begin{frame}{Lecture 9 in One Map: How an Agent Runs}
\textbf{This deck answers the second question: how does it actually work, and how do I start?}

\LecNineMap{2}{a}

\vspace{0.1cm}
\footnotesize\textit{Route: the three physical components and the tool-call loop $\to$ install and initialize $\to$ the shape of the CLI $\to$ your first session (on the fellows roster) $\to$ tokens, the context window, and the pillars that keep you honest.}
\end{frame}

% ---- T2b-F01 · MERGE-PENDING(T2-F01+T2b-F01) · TIER: CORE · from T2b:36
\begin{frame}{Topic 9.2b Agenda}
\textbf{T2 opened the hood on one agent. This interlude answers the six questions students actually ask once they have run a session --- and corrects three beliefs that are widespread and wrong.}

\vspace{0.2cm}

\begin{enumerate}
  \item \textbf{One agent, written out} --- what a ``session'' physically is, and the loop in twenty lines of Python
  \item \textbf{Three backgrounds, three agents} --- the eight steps that turn a job description into a file you can dispatch
  \item \textbf{Many agents} --- how you create a \emph{set} of them, and what actually limits the number
  \item \textbf{Who decides to dispatch?} --- the router myth, and what really happens
  \item \textbf{Agent vs.\ skill} --- a distinction most write-ups get backwards
  \item \textbf{Synthesis} --- and a lab where you design and dispatch your own agents
\end{enumerate}

\vspace{0.18cm}

\begin{takeawaybox}
\footnotesize\textbf{Why a separate deck.} T1b surveyed the \emph{landscape} of multi-agent systems; T4 taught you to \emph{use} sub-agents. Neither showed the machinery, and neither had you \emph{build} one. Everything here is checkable on your own machine --- and the lab makes you check it.
\end{takeawaybox}

\vspace{0.1cm}
{\scriptsize \textit{Reading convention, as in T1b: \textbf{[solid]} $=$ primary or citable; \textbf{[hype]} $=$ vendor self-report.}}
\end{frame}


%=============================================================================
\section{Three Components and One JSON Line}
%===========================================================

\sectiondivider{Three Components and One JSON Line}{Harness, model, files --- and the loop that joins them}

% ---- T2-F02 · KEEP · TIER: READ · from T2:46
\begin{frame}{Three Components: Harness, LLM, Your Files}

A terminal agent has three physically distinct parts --- the same anatomy in Claude Code, Codex CLI, Gemini CLI, or Aider:

\vspace{0.2cm}

\begin{columns}[T]
\begin{column}{0.55\textwidth}
\begin{enumerate}
    \item \textbf{The Harness} (lives on your machine)\\
    \small CLI software you install. Reads your files, runs shell commands, manages the conversation loop.

    \vspace{0.2cm}

    \item \textbf{The LLM} (lives on the provider's servers)\\
    \small Every message is bundled with conversation history and sent via HTTPS to the API. The ``thinking'' happens in the cloud.

    \vspace{0.2cm}

    \item \textbf{Your Files} (live on your machine)\\
    \small Python scripts, data CSVs, LaTeX---all on your disk. Code executes on your hardware with your packages.
\end{enumerate}
\end{column}
\begin{column}{0.4\textwidth}
\begin{center}
\begin{tikzpicture}[
    macbox/.style={rectangle, draw=blue!60, fill=blue!8, rounded corners,
                   minimum width=3.2cm, minimum height=1.0cm,
                   text width=3.0cm, align=center, font=\small},
    cloudbox/.style={rectangle, draw=orange!70, fill=orange!10, rounded corners,
                     minimum width=3.2cm, minimum height=1.0cm,
                     text width=3.0cm, align=center, font=\small},
    arrow/.style={-{Stealth[length=2mm]}, thick, gray}
]
    \node[macbox] (harness) {Harness + Files\\{\footnotesize Your machine}};
    \node[cloudbox, below=1.4cm of harness] (llm) {LLM\\{\footnotesize Anthropic API}};
    \draw[arrow, bend left=15] (harness) to node[right,font=\tiny] {prompt + context} (llm);
    \draw[arrow, bend left=15] (llm) to node[left,font=\tiny] {tool calls + text} (harness);
\end{tikzpicture}
\end{center}
\vspace{0.15cm}
\footnotesize
\textbf{Key:} code can run locally and many project files stay on your machine. That does \textit{not} mean the whole workflow is automatically private.
\end{column}
\end{columns}
\end{frame}

% ---- T2-F03 · MERGE-PENDING(T2-F03+T2b-F02) · TIER: CORE · from T2:93
% TODO(Phase B): One frame -- T2b-F02's three-node diagram above T2-F03's four bullets.
\begin{frame}[fragile]{The Loop in Machine Language: JSON Tool Calls}
\textbf{Lec07 T4 explained that an LLM can emit structured JSON via schema-constrained decoding. In a terminal agent, every tool call uses that mechanism.}

\textbf{Example.} When the agent decides to read a file, it emits:
\begin{lstlisting}[style=pythonstyle,language={},basicstyle=\ttfamily\scriptsize,numbers=none]
{"tool": "Read", "args": {"file_path": "fellows_seed.csv"}}
\end{lstlisting}

The harness parses the JSON, executes \texttt{Read}, and feeds the file content back as the next message; the agent reads it and decides the next action.

\begin{itemize}
    \item \textbf{Loop.} Model output (JSON) $\rightarrow$ harness execution $\rightarrow$ tool result $\rightarrow$ next model turn. Continues until the model emits a ``done'' signal or hits a stop condition.
    \item \textbf{Schema is the contract.} Each tool advertises an argument schema; outputs that fail the schema are rejected.
    \item \textbf{Same brain.} There is no separate ``planner'' --- the same transformer that writes code also decides to call tools.
    \item \textbf{Key implication.} Every action the agent takes is a \emph{next-token prediction} constrained by the available tool schema.
\end{itemize}
\end{frame}

% ---- T2b-F02 · MERGE-PENDING(T2-F03+T2b-F02) · TIER: CORE · from T2b:66
\begin{frame}{From One JSON Line to a Running Loop}
\textbf{When an agent wants to read a file, it does not read the file. It \emph{writes a request} --- one line of JSON --- and something else does the reading.}

\vspace{0.1cm}

\begin{center}
\begin{tikzpicture}[font=\scriptsize, yscale=0.78,
    b/.style={rectangle, draw=black!60, fill=VeryLightGray, rounded corners,
              minimum width=2.6cm, minimum height=0.95cm, text width=2.45cm, align=center},
    m/.style={rectangle, draw=green!45!black, fill=green!6, rounded corners,
              minimum width=2.6cm, minimum height=0.95cm, text width=2.45cm, align=center},
    arr/.style={-{Stealth[length=1.9mm]}, thick, gray!70}]
    \node[m] (llm)  at (0,0)    {\textbf{The model}\\ (on Anthropic's servers)};
    \node[b] (harn) at (4.6,0)  {\textbf{The harness}\\ Claude Code, on your laptop};
    \node[b] (disk) at (9.2,0)  {\textbf{Your disk}\\ \texttt{fellows\_seed.csv}};
    \draw[arr] (llm) -- node[above, font=\tiny, align=center] {\textbf{writes} this JSON} (harn);
    \draw[arr] (harn) -- node[above, font=\tiny] {\textbf{runs} it} (disk);
    \draw[arr] (disk) to[bend left=28] node[below, font=\tiny] {file contents go back to the model} (llm);
\end{tikzpicture}
\end{center}

\vspace{0.05cm}

\begin{columns}[T]
\begin{column}{0.47\textwidth}
\footnotesize\textbf{What the model emits} (the line above):

\vspace{0.08cm}
{\ttfamily\scriptsize \{"tool": "Read",\\ \phantom{x}"args": \{"file\_path": "\ldots"\}\}}

\vspace{0.1cm}
{\scriptsize It is \emph{text the model generated}, in a shape you told it to use. The model cannot touch your disk; the harness can.}
\end{column}
\begin{column}{0.5\textwidth}
\begin{conceptbox}
\footnotesize
\textbf{The naive mental model.} ``I open a session, the model remembers me, and somewhere a dispatcher assigns work to specialists.''

\vspace{0.12cm}
Every clause is wrong. This deck replaces it --- and the cost model, the context limits, and the sub-agent behaviour of T3b and T5b--T5c all follow.
\end{conceptbox}
\end{column}
\end{columns}

\vspace{0.08cm}
{\footnotesize\textbf{Answered next:} if nothing is stored, what is a ``session'' --- and where is ``the agent''?}
\end{frame}

% ---- T2-F04 · KEEP · TIER: READ · from T2:111
\begin{frame}{Privacy Boundary: Local Execution Is Not the Same as Local Processing}
\begin{itemize}
    \item code may execute locally on your machine
    \item many project files may stay on disk locally
    \item if file contents, prompts, or tool outputs are sent back in the conversation, those contents may still traverse the network
    \item only a fully local or institutionally private deployment changes that boundary
\end{itemize}

\vspace{0.25cm}
\textbf{Practical rule:} running code locally does not make it safe to put restricted data into the conversation --- treat everything you type or attach as leaving your machine.
\end{frame}


%=============================================================================
\section{The Session Is a List}
%===========================================================

\sectiondivider{The Session Is a List}{What a session is, and the loop that drives it}

% ---- T2b-F03 · KEEP · TIER: READ · from T2b:114
\begin{frame}{First, the Vocabulary: Who Is Speaking?}
\textbf{Every message carries a \texttt{role} label saying who produced it. There are three things you will see, and one is misleadingly named.}

\vspace{0.1cm}

\begin{center}
\footnotesize
\renewcommand{\arraystretch}{1.25}
\begin{tabular}{|p{2.4cm}|p{3.4cm}|p{4.8cm}|}
\hline
\textbf{Label} & \textbf{Who actually wrote it} & \textbf{Example} \\ \hline
\texttt{"user"} & \textbf{you} & ``Solve the model and report $r$.'' \\ \hline
\texttt{"assistant"} & \textbf{the LLM} --- Claude itself & the JSON tool request on the last slide \\ \hline
\texttt{"user"}, containing a \texttt{tool\_result} & \textbf{your harness}, reporting what a tool returned & the contents of \texttt{fellows\_seed.csv} \\ \hline
\end{tabular}
\end{center}

\vspace{0.1cm}

\begin{cautionbox}
\footnotesize\textbf{The naming trap.} ``Assistant'' is simply the API's word for \emph{the model's own turn} --- it is not a separate helper program. And tool results are labelled \texttt{"user"} even though you did not type them: the API has only two speakers, so anything sent \emph{to} the model is labelled \texttt{"user"}, whoever produced it. Read \texttt{"user"} as ``input to the model'', not as ``the human''.
\end{cautionbox}

\vspace{0.08cm}
{\footnotesize\textbf{Why it matters:} \texttt{assistant} means \emph{the model}; tool output re-enters as \texttt{user}. That is all the vocabulary this deck needs.}
\end{frame}

% ---- T2b-F04 · KEEP · TIER: CORE · from T2b:141
\begin{frame}{The Session Is a List, Not a Connection}
\textbf{The API is stateless. Each turn is a \emph{separate} request that carries the entire conversation so far --- the model remembers nothing between calls.}

\vspace{0.12cm}

\begin{center}
\begin{tikzpicture}[
    msg/.style={rectangle, draw=black!55, fill=VeryLightGray, rounded corners,
                minimum width=2.5cm, minimum height=0.42cm, font=\tiny, align=center},
    new/.style={rectangle, draw=RedTitle!75, fill=blue!10, rounded corners,
                minimum width=2.5cm, minimum height=0.42cm, font=\tiny, align=center},
    hd/.style={font=\scriptsize\bfseries},
    arr/.style={-{Stealth[length=1.8mm]}, thick, gray!65}
]
    \node[hd] at (0,1.15) {Turn 1};
    \node[new] (a1) at (0,0.6) {user: ``solve it''};
    \node[draw=gray!45, dashed, fit=(a1), inner sep=2pt] {};

    \node[hd] at (3.6,1.15) {Turn 2};
    \node[msg] (b1) at (3.6,0.6) {user: ``solve it''};
    \node[msg] (b2) at (3.6,0.12) {assistant: tool\_use};
    \node[new] (b3) at (3.6,-0.36) {user: tool\_result};
    \node[draw=gray!45, dashed, fit=(b1)(b3), inner sep=2pt] {};

    \node[hd] at (7.2,1.15) {Turn 3};
    \node[msg] (c1) at (7.2,0.6) {user: ``solve it''};
    \node[msg] (c2) at (7.2,0.12) {assistant: tool\_use};
    \node[msg] (c3) at (7.2,-0.36) {user: tool\_result};
    \node[msg] (c4) at (7.2,-0.84) {assistant: tool\_use};
    \node[new] (c5) at (7.2,-1.32) {user: tool\_result};
    \node[draw=gray!45, dashed, fit=(c1)(c5), inner sep=2pt] {};

    \draw[arr] (1.35,0.1) -- (2.25,0.1);
    \draw[arr] (4.95,-0.1) -- (5.85,-0.1);

    \node[font=\tiny, gray!75, align=center] at (1.8,-0.35) {new request,\\\emph{whole list again}};
    \node[font=\tiny, gray!75, align=center] at (5.4,-0.55) {new request,\\\emph{whole list again}};

    \node[font=\scriptsize, align=center] at (10.2,0.0)
        {\textbf{Blue $=$ the only}\\\textbf{new bytes}\\[3pt]
         Grey $=$ re-sent\\verbatim, and\\re-charged};
\end{tikzpicture}
\end{center}

\vspace{0.05cm}

\begin{columns}[T]
\begin{column}{0.52\textwidth}
\begin{itemize}\footnotesize
    \item \textbf{``Re-send everything'' is literal:} turn 3 is a fresh request carrying turns 1--3 in full
    \item No open connection, no server-side memory of you
    \item ``Context window full'' means \emph{this list} outgrew the limit
    \item You are billed for the whole list each turn
\end{itemize}
\end{column}
\begin{column}{0.46\textwidth}
\begin{takeawaybox}
\footnotesize\textbf{Where the list starts and stops.} \texttt{/clear} does exactly one thing: it throws the list away and starts an empty one. Compaction replaces the list with a summary of itself. Both are list operations --- there is no ``connection'' to reset.
\end{takeawaybox}
\end{column}
\end{columns}
\end{frame}

% ---- T2b-F05 · EDIT · TIER: READ · from T2b:204
\begin{frame}{Four Moving Parts, One Function Call}
\textbf{An agent is not a service you connect to. It is four arguments you assemble and send.}

\vspace{0.12cm}

\begin{center}
\footnotesize
\renewcommand{\arraystretch}{1.3}
\begin{tabular}{|p{1.9cm}|p{3.4cm}|p{5.4cm}|}
\hline
\textbf{Argument} & \textbf{What it is} & \textbf{Concretely} \\ \hline
\texttt{model} & \textbf{which LLM} to send to & the string \texttt{"claude-opus-5"} --- just a name; you are picking a brain \\ \hline
\texttt{system} & the \textbf{standing instructions}: role, rules, style, read \emph{before} the conversation & ``You are a computational macroeconomist. Never invent data.'' This is what \texttt{CLAUDE.md} becomes (T2b, T5b) \\ \hline
\texttt{tools} & the \textbf{menu of actions} --- name, description, argument schema. \emph{Descriptions only; no code is sent} & ``\texttt{Read}: read a file. Arguments: \texttt{file\_path} (string).'' The model picks from this menu \\ \hline
\texttt{messages} & the \textbf{conversation so far} & every \texttt{user} / \texttt{assistant} / \texttt{tool\_result} turn, in order \\ \hline
\end{tabular}
\end{center}

\vspace{0.06cm}
{\footnotesize\textbf{What comes back:} ordinary text, or \texttt{tool\_use} blocks --- the model asking you to run something from the menu. Change \texttt{system} and \texttt{tools} and you have a \emph{different agent}; there is no ``agent object'' to create or shut down.}
\end{frame}

% ---- T2b-F06 · KEEP · TIER: LAB · from T2b:226
\begin{frame}[fragile]{The Loop in Twenty Lines --- Written by You, Once}
\textbf{You write this loop once. The model never writes or runs it --- it only supplies the replies the loop reads.}

\vspace{0.04cm}

\begin{lstlisting}[style=pythonstyle,basicstyle=\ttfamily\fontsize{5.4}{6.2}\selectfont]
messages = [{"role": "user", "content": "Solve the model and report r."}]

while True:
    resp = client.messages.create(
        model="claude-opus-5",
        max_tokens=16000,
        system="You are a computational macroeconomist.",
        tools=TOOLS,                      # schemas, not functions
        messages=messages,                # the whole list, every time
    )
    if resp.stop_reason != "tool_use":    # model is done talking
        break
    messages.append({"role": "assistant", "content": resp.content})

    results = []
    for b in resp.content:                # may be SEVERAL tool calls
        if b.type == "tool_use":
            results.append({"type": "tool_result",
                            "tool_use_id": b.id,      # must match
                            "content": run_tool(b.name, b.input)})
    messages.append({"role": "user", "content": results})
\end{lstlisting}

\vspace{0.02cm}
{\scriptsize\textbf{Not automatic:} this file --- a person wrote it. \textbf{Automatic:} \emph{which} tool, \emph{which} arguments, when to stop --- chosen by the model inside \texttt{create()}, line 4. \textbf{Line 9 with line 21:} the list only grows and is re-sent whole every pass.}
\end{frame}

% ---- T2b-F07 · KEEP · TIER: READ · from T2b:259
\begin{frame}{ReAct, Named --- and Where the Textbook Diagram Is Now Dated}
\textbf{The loop you just read has a name in the literature: ReAct --- reason, act, observe, repeat \textnormal{(Yao et al., 2022)}.}

\vspace{0.1cm}

\begin{columns}[T]
\begin{column}{0.5\textwidth}
\begin{center}
\begin{tikzpicture}[
    st/.style={rectangle, draw=RedTitle!65, fill=blue!6, rounded corners,
               minimum width=2.3cm, minimum height=0.62cm, font=\scriptsize, align=center},
    arr/.style={-{Stealth[length=1.9mm]}, thick, gray!70}
]
    \node[st] (r) at (0,1.5)    {\textbf{Reason}};
    \node[st] (a) at (2.1,0.0)  {\textbf{Act}};
    \node[st] (o) at (0,-1.5)   {\textbf{Observe}};
    \node[st] (d) at (-2.1,0.0) {\textbf{Done?}};

    \draw[arr] (r) to[bend left=22] node[font=\tiny, above right=-1pt] {line 4} (a);
    \draw[arr] (a) to[bend left=22] node[font=\tiny, below right=-1pt] {line 20} (o);
    \draw[arr] (o) to[bend left=22] node[font=\tiny, below left=-1pt] {line 21} (d);
    \draw[arr] (d) to[bend left=22] node[font=\tiny, above left=-1pt] {line 11} (r);
\end{tikzpicture}
\end{center}
\end{column}
\begin{column}{0.48\textwidth}
\begin{itemize}\footnotesize
    \item \textbf{Reason} $=$ the model's turn (line 4)
    \item \textbf{Act} $=$ your harness runs the tool (line 20)
    \item \textbf{Observe} $=$ the \texttt{tool\_result} goes back (line 21)
    \item \textbf{Done?} $=$ the \texttt{stop\_reason} test (line 11)
\end{itemize}
\end{column}
\end{columns}

\vspace{0.15cm}

\begin{cautionbox}
\footnotesize\textbf{Where the 2022 picture is dated.} The original ReAct prompted the model to emit literal \texttt{Thought:} / \texttt{Action:} text that a parser scraped. Current models emit \emph{structured} \texttt{tool\_use} blocks and interleave reasoning with tool calls rather than strictly alternating. \textbf{Treat ReAct as the ancestor of the loop, not as its wire format} --- if you implement the 2022 string-parsing version today you are rebuilding something the API already gives you typed.
\end{cautionbox}
\end{frame}

% ---- T2b-F08 · EDIT · TIER: READ · from T2b:301
\begin{frame}{So Where Does the Agent Live?}
\textbf{The JSON frame said: \emph{``the agent reads it and decides the next action.''} That is convenient shorthand. Taken literally it is misleading --- no single place contains the agent.}

\vspace{0.1cm}

\begin{center}
\begin{tikzpicture}[font=\scriptsize, yscale=0.92,
    loc/.style={rectangle, draw=black!60, fill=VeryLightGray, rounded corners,
                minimum width=3.5cm, minimum height=1.75cm, text width=3.35cm, align=center},
    hd/.style={font=\scriptsize\bfseries}]
    \node[hd] at (0,1.15)   {Your laptop};
    \node[loc] at (0,0)     {the \textbf{harness}: the loop\\ the \textbf{tool code}\\ \textbf{the list}, in memory};
    \node[hd] at (4.5,1.15) {The wire};
    \node[loc, draw=RedTitle!60, fill=blue!5] at (4.5,0)
        {one \textbf{HTTPS request}\\ per turn, carrying\\ the \emph{whole} list};
    \node[hd] at (9.0,1.15) {Anthropic's servers};
    \node[loc, draw=green!45!black, fill=green!5] at (9.0,0)
        {the \textbf{model weights}\\ \emph{stateless} --- keeps\\ nothing between calls};
    \draw[-{Stealth[length=1.8mm]}, thick, gray!70] (1.85,0.3) -- (2.7,0.3);
    \draw[-{Stealth[length=1.8mm]}, thick, gray!70] (6.3,-0.3) -- (7.2,-0.3);
\end{tikzpicture}
\end{center}

\vspace{0.1cm}

\begin{takeawaybox}
\footnotesize\textbf{The agent is the loop, not a location.} ``The agent'' names a \emph{behaviour} that appears when a stateless model is driven by a stateful loop. Delete the loop and there is no agent --- only a model that answers one question. Read T2's sentence as: \emph{the harness appended the tool result to the list, and the next request carried that list to the model.}
\end{takeawaybox}

\vspace{0.08cm}
{\footnotesize\textbf{Why this is the practical frame:} it tells you where each cost and limit lives --- tokens and the context window on the \emph{wire and the server}, tool permissions and your data on \emph{your laptop} (T2's privacy boundary), and the memory in a list \emph{you} own and can clear.}
\end{frame}


%=============================================================================
\section{Worked Example --- GDP Growth by Party}
%===========================================================

\sectiondivider{Worked Example --- GDP Growth by Party}{One real research question, traced message by message}

% ---- T2b-F33 · KEEP · TIER: CORE · from T2b:1371
% TODO(Phase B): Review A4: footnote the recorded transcript (date, harness + version, model) once it exists in labs/Lec09_Agent_Lab/transcripts/.
\begin{frame}{The Task --- and the Trap Inside It}
\textbf{One prompt, typed once:}

\vspace{0.12cm}
\begin{conceptbox}
\footnotesize\textit{``Get US real GDP growth since 2000. Get which party held the White House over the same period. Then plot and \textbf{rationalize how party affiliation affects GDP growth}.''}
\end{conceptbox}

\vspace{0.15cm}

\begin{columns}[T]
\begin{column}{0.49\textwidth}
\footnotesize\textbf{Three jobs, not one}
\begin{enumerate}\footnotesize
    \item fetch a macro series (a \emph{fact})
    \item build a party timeline (a \emph{fact})
    \item plot it (\emph{mechanical})
    \item explain the relationship (\emph{a causal claim})
\end{enumerate}

\vspace{0.1cm}
{\scriptsize Steps 1--3 an agent can do. Step 4 is where the research is --- and where the prompt is booby-trapped.}
\end{column}
\begin{column}{0.49\textwidth}
\begin{cautionbox}
\footnotesize\textbf{``Rationalize'' is a leading verb.} It presupposes an effect and asks for a justification. A compliant model will supply one. Watch, in the trace that follows, the exact turn where a defensible data task turns into an indefensible causal claim --- and notice that nothing in the machinery stops it.
\end{cautionbox}
\end{column}
\end{columns}

\vspace{0.12cm}
\footnotesize\textbf{A chat model with no tools fails this instantly and invisibly}: it will recite GDP numbers from memory. They will look right. Some will be wrong, and none will be checkable.
\end{frame}

% ---- T2b-F34 · EDIT · TIER: CORE · from T2b:1405
% TODO(Phase B): Review A4: sync the five requests with the recorded run; if the real run differs, the slide changes. Review A2: label the server "LLM (vendor)" rather than one company.
\begin{frame}{One Agent: Every Message That Crosses the Wire}
\textbf{Five requests. Each carries the \emph{entire} list built so far.}

\vspace{0.02cm}

\begin{center}
\begin{tikzpicture}[font=\tiny, yscale=0.80,
    ln/.style={draw=black!35, dashed},
    hd/.style={font=\scriptsize\bfseries, align=center},
    a/.style={-{Stealth[length=1.6mm]}, thick},
    loc/.style={a, gray!75},
    wire/.style={a, RedTitle!75}]
    \node[hd] (u) at (0,0.55)   {You};
    \node[hd] (h) at (4.3,0.55) {Harness\\ \textnormal{\tiny(your laptop)}};
    \node[hd] (l) at (9.6,0.55) {LLM\\ \textnormal{\tiny(Anthropic)}};
    \draw[ln] (0,0.1) -- (0,-6.5);
    \draw[ln] (4.3,0.1) -- (4.3,-6.5);
    \draw[ln] (9.6,0.1) -- (9.6,-6.5);

    \draw[loc] (0,-0.25) -- node[above]{the prompt} (4.3,-0.25);
    \draw[wire] (4.3,-0.7) -- node[above]{\textbf{req 1} = system + tools + [prompt]} (9.6,-0.7);
    \draw[wire] (9.6,-1.15) -- node[above]{\texttt{tool\_use}: Bash \emph{fetch GDP series}} (4.3,-1.15);
    \draw[loc] (4.3,-1.6) -- node[above, xshift=2mm]{runs it locally} (3.0,-1.6);
    \draw[wire] (4.3,-2.05) -- node[above]{\textbf{req 2} = list + \texttt{tool\_result}(100 rows)} (9.6,-2.05);
    \draw[wire, orange!80!black] (9.6,-2.5) -- node[above, orange!70!black]{\texttt{tool\_use}: Write \emph{party\_by\_year.csv}} (4.3,-2.5);
    \node[font=\tiny, orange!70!black, anchor=west] at (9.75,-2.5) {\textbf{no fetch}};
    \node[font=\tiny, orange!70!black, anchor=west, align=left] at (9.75,-2.9) {from the model's\\ own memory};
    \draw[wire] (4.3,-2.95) -- node[above]{\textbf{req 3} = list + \texttt{tool\_result}(ok)} (9.6,-2.95);
    \draw[wire] (9.6,-3.4) -- node[above]{\texttt{tool\_use}: Bash \emph{python plot.py}} (4.3,-3.4);
    \draw[loc, red!60!black] (4.3,-3.85) -- node[above, red!55!black, xshift=1mm]{\textbf{error}: no display backend} (2.4,-3.85);
    \draw[wire, red!60!black] (4.3,-4.3) -- node[above]{\textbf{req 4} = list + \texttt{tool\_result}(\emph{is\_error})} (9.6,-4.3);
    \draw[wire] (9.6,-4.75) -- node[above]{\texttt{tool\_use}: Bash \emph{retry with Agg}} (4.3,-4.75);
    \draw[wire] (4.3,-5.2) -- node[above]{\textbf{req 5} = list + \texttt{tool\_result}(figure written)} (9.6,-5.2);
    \draw[wire] (9.6,-5.65) -- node[above]{text + \texttt{stop\_reason}: \texttt{end\_turn}} (4.3,-5.65);
    \draw[loc] (4.3,-6.1) -- node[above]{the answer, and a \texttt{.png}} (0,-6.1);
\end{tikzpicture}
\end{center}

\vspace{0.02cm}
{\scriptsize\textbf{Blue crosses the network and is billed; grey stays on your laptop, free.} The list grew to eleven messages --- \textbf{req 5} re-sent all of them.}
\end{frame}

% ---- T2b-F35 · EDIT · TIER: READ · from T2b:1447
\begin{frame}{What Went Into \textnormal{\texttt{req 1}} --- and Never Changed}
\textbf{This is T2b's \emph{initialization}, seen from the wire: those five steps exist to build exactly two arguments --- then re-sent byte-identical on every later request.}

\vspace{0.1cm}

\begin{center}
\footnotesize
\renewcommand{\arraystretch}{1.2}
\begin{tabular}{|p{1.5cm}|p{4.6cm}|p{5.0cm}|}
\hline
& \textbf{Assembled from} & \textbf{In this example} \\ \hline
\texttt{system} & T2b's init steps 1--4: global settings, \texttt{CLAUDE.md}, \texttt{.claude/rules/*.md} & ``You are Claude Code. You have these tools\ldots'' $+$ whatever house rules T5b has you write \\ \hline
\texttt{tools} & the schema of every tool the harness is willing to run & \texttt{Bash}, \texttt{Read}, \texttt{Write}, \texttt{Grep}\ldots\ names, descriptions, argument types --- \emph{no code} \\ \hline
\end{tabular}
\end{center}

\vspace{0.1cm}

\begin{takeawaybox}
\footnotesize\textbf{Yes --- the model writes the code and asks your machine to run it.} A \texttt{Bash} call is the model generating the \emph{text} of a shell command as an argument. It executes nothing; your harness runs that string on your laptop, with your files and permissions.
\end{takeawaybox}

\vspace{0.06cm}

\begin{columns}[T]
\begin{column}{0.47\textwidth}
{\footnotesize\textbf{The safety boundary is not the model's judgement} --- it is T2b's permission dial, between the generated string and your shell.}
\end{column}
\begin{column}{0.51\textwidth}
{\footnotesize\textbf{\texttt{system} is the only text still guaranteed to be there on turn 200} --- everything else drifts or compacts away. That is the mechanical reason \texttt{CLAUDE.md} earns its keep.}
\end{column}
\end{columns}
\end{frame}

% ---- T2b-F36 · EDIT · TIER: READ · from T2b:1481
\begin{frame}{Reading the Trace: Four Things Just Happened}
\textbf{Every claim from the session section shows up in that one picture.}

\vspace{0.15cm}

\begin{center}
\footnotesize
\renewcommand{\arraystretch}{1.28}
\begin{tabular}{|p{2.6cm}|p{4.1cm}|p{4.5cm}|}
\hline
\textbf{Moment} & \textbf{What it looked like} & \textbf{What it proves} \\ \hline
\textbf{The loop ran itself} & five \texttt{tool\_use} $\rightarrow$ \texttt{tool\_result} round trips, no human input & nobody scheduled these steps; the model chose each one (above) \\ \hline
\textbf{The error was recovered} & backend error came back as a \texttt{tool\_result} with \texttt{is\_error}, and the next turn retried & this is ``Observe'' in ReAct, and why a failed tool must still \emph{return} (above) \\ \hline
\textbf{Cost grew quadratically} & req 5 re-sent everything from req 1 & the list is the cost driver, not the number of questions (above) \\ \hline
\textbf{One step had no tool} & the party timeline was \emph{written from memory} (orange, previous slide) & the weakest link: in the transcript it looks identical to a fetched fact \\ \hline
\end{tabular}
\end{center}

\vspace{0.1cm}

\begin{cautionbox}
\footnotesize\textbf{Where a referee would stop you.} The GDP series has provenance --- a URL, a vintage, a retrieval date. The party timeline has none; it is model recall typed into a CSV. \textbf{Ask of every number an agent hands you: which tool produced this, and can I re-run it?}
\end{cautionbox}
\end{frame}

% ---- T2b-F37 · KEEP · TIER: READ · from T2b:1506
\begin{frame}{The Same Task with Three Agents}
\textbf{The useful question is not ``is this faster'' but \emph{which steps are independent, and which one needs an adversary}.}

\vspace{0.05cm}

\begin{center}
\begin{tikzpicture}[font=\tiny, yscale=0.80,
    o/.style={rectangle, draw=RedTitle!70, fill=blue!7, rounded corners,
              minimum width=3.0cm, minimum height=0.7cm, align=center},
    w/.style={rectangle, draw=green!45!black, fill=green!6, rounded corners,
              minimum width=2.9cm, minimum height=1.0cm, text width=2.75cm, align=center},
    r/.style={rectangle, draw=red!55!black, fill=red!5, rounded corners,
              minimum width=2.9cm, minimum height=1.0cm, text width=2.75cm, align=center},
    a/.style={-{Stealth[length=1.7mm]}, thick, gray!70}]
    \node[o] (orc) at (5.2,1.5) {\textbf{Orchestrator} --- owns the plot and the write-up};
    \node[w] (A) at (0.6,-0.3) {\textbf{data-fetcher}\\ GDP series $+$ URL,\\ vintage, date};
    \node[w] (B) at (5.2,-0.3) {\textbf{timeline-builder}\\ party by year,\\ \emph{with a citation}};
    \node[r] (C) at (9.8,-0.3) {\textbf{domain-reviewer}\\ attack the causal\\ claim};
    \draw[a] (orc) -- (A);  \draw[a] (orc) -- (B);  \draw[a] (orc) -- (C);
    \node[font=\tiny, gray!70] at (2.9,-1.35) {parallel: independent facts};
    \node[font=\tiny, red!55!black] at (9.8,-1.35) {sequential: needs the draft first};
\end{tikzpicture}
\end{center}

\vspace{0.12cm}

\begin{columns}[T]
\begin{column}{0.49\textwidth}
\footnotesize\textbf{What fan-out actually buys}
\begin{itemize}\footnotesize
    \item the two \emph{facts} are independent $\rightarrow$ genuinely parallel
    \item 100 rows land in \texttt{data-fetcher}'s context, \textbf{not} the orchestrator's
    \item one job each, so ``cite your source'' is enforceable
\end{itemize}
\end{column}
\begin{column}{0.49\textwidth}
\footnotesize\textbf{What it does \emph{not} buy}
\begin{itemize}\footnotesize
    \item the plot is one small step --- splitting it adds turns, not speed
    \item three agents cost more tokens for the \emph{same} numbers
    \item nothing here makes the causal claim more true
\end{itemize}
\end{column}
\end{columns}

\vspace{0.04cm}
{\footnotesize\textbf{Only \texttt{domain-reviewer} changes the answer} --- and it cannot run in parallel: it needs something to attack.}
\end{frame}

% ---- T2b-F38 · EDIT · TIER: CORE · from T2b:1555
\begin{frame}{What the Reviewer Says --- and What No Agent Can Fix}
\textbf{Given the draft, a competent \texttt{domain-reviewer} returns roughly this:}

\vspace{0.1cm}

\begin{cautionbox}
\footnotesize
``The correlation is real but you cannot read it as an effect. Since 2000 you have about \textbf{six administrations} --- your effective $n$ is single digits, not the 100 quarters you plotted. Presidents inherit economies, business cycles are not assigned at random, and Congress is unmodelled. \textbf{Blinder and Watson (2016)} find a large gap over 1949--2012 (about \textbf{1.79pp} more annual real GDP growth under Democratic presidents, 16 terms) and attribute most of it to \emph{oil shocks, productivity shocks, and foreign demand} --- not to policy. Report the correlation, state the identification problem, and drop the word `affects'.'' \textbf{[solid]}
\end{cautionbox}

\vspace{0.13cm}

\begin{columns}[T]
\begin{column}{0.5\textwidth}
\begin{takeawaybox}
\footnotesize\textbf{This is the case for multi-agent, stated honestly.} Not speed. A second agent whose job is to \emph{disagree} caught a claim the first agent was happy to make --- because you told it to.
\end{takeawaybox}
\end{column}
\begin{column}{0.48\textwidth}
\begin{cautionbox}
\footnotesize\textbf{And the limit.} Every agent here inherited your word \emph{``rationalize''}. A more agreeable reviewer would have obliged. \textbf{No number of agents fixes a leading prompt} --- that is T5c's Context Dilemma, in one word of your own typing.
\end{cautionbox}
\end{column}
\end{columns}

\vspace{0.1cm}
{\footnotesize\textbf{The fix is upstream and costs nothing:} ask ``\emph{is} there a relationship, and what would identify it?'' --- then the same machinery works for you instead of against you.}
\end{frame}

% ---- NEW · TIER: READ
% TODO(Phase B): Add the 9.2 economic twin (what am I paying for?) in Phase B.
\begin{frame}{Bridge: Put the Loop on Your Own Machine}
\textbf{You have watched the loop on the wire, once with one agent and once with three.}

\vspace{0.25cm}

\begin{itemize}
    \item A session is a list the harness re-sends; the agent is the loop, not a place
    \item Every number an agent hands you should name the tool that produced it
    \item \textbf{Next (T2b):} put that loop on your own machine --- install it, learn its files and commands, run your first session
\end{itemize}
\end{frame}

%===========================================================
\end{document}
